% !TEX root = ../Monografia.tex
\chapter{Discussão e Conclusões}
\label{cap:conclusao}

Este trabalho perguntou se a estrutura de relações que um LLM extrai das notícias contém informação sobre
a covariância futura dos retornos que não está na covariância histórica nem nos modelos de setor e de
fatores, e se essa informação produz carteiras melhores fora da amostra. A resposta tem duas partes. Sim,
o grafo de notícias contém essa informação: ele prevê a correlação futura entre pares de ações além do
histórico, do setor e da co-menção. Não, essa informação não produz carteiras melhores: em nove desenhos de
carteira, nenhum estimador com notícias reduziu de forma significativa o risco realizado em relação à sua
referência sem notícias, e em dois desenhos a carteira com notícias foi significativamente pior.

%--------------------------------------------------------------------
\section{Síntese dos resultados}
\label{sec:sintese}

\paragraph{A base de relações.} O Qwen2.5-14B-Instruct leu 572.762 artigos do FNSPID, cerca de 36\% do
corpus elegível, e extraiu 1.270.582 relações, das quais 86,1\% passaram pelo filtro de ruído. Depois da
ligação aos \emph{tickers}, elas formam um grafo dinâmico com 255.384 nós e 1.064.801 eventos entre 2011 e
2023. Na anotação automática, cerca de quatro em cada cinco tuplas descrevem uma relação que o texto de
fato afirma, e o sinal está certo em mais de 90\% das relações existentes.

\paragraph{H1: as notícias preveem a correlação.} Na correlação residual ao fator de mercado, os atributos
do grafo reduzem o erro de previsão fora da amostra em relação à persistência e ao setor em 84\% dos meses,
com estatística $t$ de 8,8, e a ampliação da amostra de artigos deixou o ganho igual ou um pouco maior. Os
dois canais previstos pela teoria aparecem nos dados. A exposição comum a entidades não negociadas é o
mais forte e sobrevive aos controles de setor e de atenção da mídia num universo que não foi escolhido
pelas notícias (H1a). O tipo de relação direta importa, com a concorrência à frente (H1b). Em magnitude, o
efeito é pequeno: cerca de três décimos de ponto de $R^2$ e quatro milésimos de correlação por
desvio-padrão de exposição comum.

\paragraph{H2: as notícias não melhoram as carteiras.} A Tabela~\ref{tab:sintese-h2} mostra o padrão. Sem
restrição de peso, as notícias empatam com a referência setorial ou a pioram (D1, D2 e D9); com
restrições, empatam (D3 a D7); no mínimo \emph{tracking error}, pioram de forma crescente com o número de
ações (D8). O resultado se repete em universos de 50, 100 e 200 ações e nas duas amostras de artigos, e,
sempre que as notícias pareceram ajudar, o placebo ajudou tanto quanto elas.

Os dois resultados não se contradizem. O grafo sabe algo sobre o co-movimento, mas o que ele sabe é
pouco, está na média dos pares e já está, em boa parte, no setor e nos fatores. A carteira de variância
mínima, por sua vez, não depende da precisão média das correlações, e sim da estrutura da matriz que o
otimizador inverte.

%--------------------------------------------------------------------
\section{Por que a informação não vira carteira melhor}
\label{sec:mecanismos}

\paragraph{O ganho de previsão é pequeno e está no lugar errado.} O ganho das notícias melhora a correlação
\emph{média} dos pares, enquanto a GMV depende sobretudo dos autovalores e autovetores da matriz. O D1
mostra essa diferença de forma direta: o alvo com setor e notícias teve o menor erro quadrático de
correlação entre todos os estimadores e, mesmo assim, uma carteira 3,7\% pior que a correlação constante; o
fator único de \citet{ledoit2003}, ao contrário, teve o maior erro de correlação e a melhor carteira.
Melhorar a previsão das correlações par a par não é o mesmo que melhorar a matriz para a otimização.

\paragraph{A inversão amplifica o erro da estrutura acrescentada.} As notícias acrescentam uma estrutura
esparsa e de pequena magnitude à correlação residual. No D2, essa estrutura aumentou o número de condição
da matriz em todos os universos, de 6\% com 50 ações a 45\% com 200. Como a GMV inverte a matriz, os
autovalores pequenos, os mais afetados pelo erro, recebem os maiores pesos, e o otimizador explora
justamente o erro da estimativa \citep{lopezdeprado2016}. O Black-Litterman do D9 mostra que o problema
não é só de condicionamento: nele o número de condição até caiu, mas milhares de covariâncias foram
movidas por visões cujo erro tem a mesma ordem de grandeza do sinal que carregam, e a variância realizada
subiu 14\%.

\paragraph{As restrições já regularizam a matriz.} Com venda a descoberto proibida e teto de peso, a
carteira de variância mínima equivale a usar uma matriz encolhida \citep{jagannathan2003}. No D3 e nos
otimizadores do D4 ao D7, todos os estimadores baseados em fatores ficaram a menos de um ponto de
volatilidade uns dos outros, e o placebo empatou com as notícias verdadeiras. Quando as restrições já
fazem o trabalho de estabilizar a matriz, não sobra espaço para a informação das notícias se manifestar.

\paragraph{O que o grafo sabe, o setor já sabe em boa parte.} Na regressão com todos os atributos, os
coeficientes de setor caem a menos da metade quando as notícias entram: as duas fontes explicam, em grande
parte, a mesma coisa. Cerca de metade das arestas entre ações liga empresas do mesmo setor, e o modelo de
fator de mercado com resíduo encolhido para o setor, sem nenhuma notícia, já empata com o fator único de
Ledoit e Wolf.

Em conjunto, os mecanismos explicam por que o resultado negativo de H2 não vem da falta de informação. A
informação existe e é medida com precisão (H1), mas é pequena demais e está numa dimensão da matriz, a
média das correlações par a par, que a otimização de carteiras não valoriza. A pergunta relevante para a
carteira não é se um sinal prevê correlações, mas se ele melhora a matriz nas direções em que o otimizador
a usa.

%--------------------------------------------------------------------
\section{O sinal extraído pelo LLM e a correlação}
\label{sec:sinal-llm}

O resultado mais inesperado do trabalho está na análise do sinal (Seção~\ref{sec:sinal}): a aresta direta
com sinal tem efeito \emph{negativo} sobre a correlação futura, quando o esperado era positivo. O problema
não está na qualidade da extração, pois o sinal das relações existentes está certo em mais de 90\% dos
casos. O problema é o que o sinal significa. Para o LLM, e para o texto que ele lê, uma relação negativa
descreve o efeito econômico de uma empresa sobre a outra: o lançamento de um produto por um concorrente é
ruim para a rival. Mas concorrentes do mesmo ramo estão expostos aos mesmos choques de demanda, de custo e
de regulação, e suas ações andam juntas.

O achado tem uma consequência para quem usa o sentimento extraído por LLMs como indicador de co-movimento:
\textbf{a direção de um efeito econômico não é a direção da correlação entre os retornos}. Para estimar a
covariância, a existência e a intensidade do vínculo são a informação útil; o sinal só seria útil numa
janela de evento, e não numa correlação mensal. Nas carteiras, a questão do sinal não interfere: os
atributos que alimentam os estimadores usam o módulo dos vínculos.

%--------------------------------------------------------------------
\section{Comparação com a literatura}
\label{sec:comparacao}

Os três antecedentes mais próximos relatam ganhos de carteira com informação textual
(Seção~\ref{sec:texto-sigma}). A primeira parte desses resultados se confirma aqui: as notícias preveem a
correlação, e o grafo extraído pelo LLM prevê mais que a simples co-menção de \citet{tang2019}. A segunda
não se confirma, e duas diferenças de desenho ajudam a entender por quê.

A primeira é a referência. Aqui, toda matriz com notícias foi comparada a uma referência que já usa o
setor e, em quase todos os desenhos, um fator de mercado, além de um placebo com a mesma esparsidade.
Contra referências mais fracas, as notícias também ganhariam: no D1, o alvo com notícias tem 15\% menos
variância que a covariância amostral, mas o placebo tem a mesma vantagem. Um ganho medido contra a
covariância amostral mistura o efeito do conteúdo do texto com o efeito genérico de impor alguma estrutura
à matriz.

A segunda é a janela. \citet{becquin2023} estimam a covariância com três meses de dados, um regime em que
a amostra é muito pouco informativa e qualquer \emph{prior} razoável tende a ajudar. Aqui, a janela foi
de 252 pregões, e universos maiores, com menos observações por parâmetro, desfavoreceram as notícias em
vez de favorecê-las.

O resultado converge com a literatura que mostra a dificuldade de superar estimadores simples fora da
amostra \citep{demiguel2009}: as estruturas que funcionam bem na carteira são as que reduzem o número de
parâmetros efetivos da matriz, e não as que acrescentam informação par a par.

%--------------------------------------------------------------------
\section{Limitações}
\label{sec:limitacoes}

\begin{itemize}
  \item \textbf{Amostra de artigos.} Por limite de orçamento, o LLM leu cerca de 36\% dos artigos
        elegíveis. A ampliação de 25\% para 36\% não mudou nenhuma conclusão, o que sugere, sem garantir,
        que o corpus completo também não mudaria.
  \item \textbf{Qualidade da extração.} A qualidade das tuplas foi medida por outro LLM, e não por
        anotadores humanos, em 567 tuplas, e o \emph{recall} não foi medido. A classe da relação é a
        medida mais fraca (62\% de acerto), e a concorrência, o grupo de maior efeito em H1b, é a classe
        menos confiável.
  \item \textbf{Preços e universos.} O painel de preços é formado por coletas emendadas, não cobre todo o
        S\&P~500 e perde 3.137 \emph{tickers} em 2020. Essas falhas afetam igualmente os estimadores com e
        sem notícias. O universo das carteiras é escolhido pelo número de citações, a própria variável em
        teste, e a classificação setorial é a atual, aplicada a todo o período.
  \item \textbf{Tempo e memória do modelo.} As notícias não têm hora de publicação, e o grafo descarta a
        informação do próprio dia. O período avaliado não é posterior aos dados de treino do modelo de
        linguagem; o teste de anonimização não encontrou vazamento, mas tem pouco poder estatístico.
  \item \textbf{Alcance.} Os resultados valem para ações americanas de grande capitalização, entre 2013 e
        2023, com rebalanceamento mensal e as carteiras estudadas.
\end{itemize}

%--------------------------------------------------------------------
\section{Trabalhos futuros}
\label{sec:futuros}

\begin{enumerate}
  \item \textbf{Usos que valorizam a correlação par a par.} A informação do grafo melhora a previsão da
        correlação de cada par, e não a estrutura de autovalores. Aplicações em que o par importa por si,
        como a seleção de pares para \emph{hedge} e a previsão de contágio entre empresas ligadas, podem
        aproveitar melhor esse sinal.
  \item \textbf{A exposição comum como fator.} H1a sugere que a exposição a entidades como o Federal
        Reserve ou o petróleo descreve cargas em fatores de risco. Usar o grafo para escolher fatores
        observáveis e orientar a estimação das cargas, em vez de alterar correlações residuais par a par,
        não aumenta o número de parâmetros efetivos da matriz.
  \item \textbf{Eventos e horizontes curtos.} Um estudo com notícias datadas por hora e janelas de evento
        poderia testar se o sinal do LLM prevê o co-movimento nos dias seguintes à notícia, onde a direção
        do efeito econômico deveria importar.
  \item \textbf{Extração mais completa e auditável.} Processar o corpus inteiro, guardar o trecho que
        justifica cada tupla e medir precisão e \emph{recall} com anotação humana permitiriam separar o
        que é limite da informação das notícias do que é limite da extração.
\end{enumerate}

Mostrar que um sinal extraído do texto prevê alguma grandeza financeira é o primeiro passo, e não o
último. Entre a previsão e a carteira há um otimizador que amplifica erros, restrições que já regularizam
a matriz e referências simples difíceis de superar. Um protocolo com pré-registro, placebo e confirmação
em amostras e universos diferentes é o que permite distinguir, nesse caminho, a informação real do efeito
genérico de regularizar.

\clearpage
